\section{Conservation laws and nonzero spatial modes}
\label{sec:conservation-independent}

Let $S\subseteq\mathbb R^n$ be the stoichiometric subspace and suppose
$J(\mathbb R^n)\subseteq S$, as occurs for a mass-action Jacobian
$J=\Gamma M$ with $S=\operatorname{im}\Gamma$.

\begin{proposition}[Semisimple conservation block]
\label{prop:conservation-splitting-independent}
If $J(S)\subseteq S$ and the restriction $J|_S$ is Hurwitz, then
\[
 \mathbb R^n=S\oplus\ker J,
 \qquad
 J\sim\begin{pmatrix}J|_S&0\\0&0\end{pmatrix}.
\]
In particular, all conservation-induced zero eigenvalues are semisimple and,
with $s=\dim S$,
\[
 \det(\lambda I-J)=\lambda^{n-s}
 \det(\lambda I_s-J|_S)>0\qquad(\lambda>0).
\]
\end{proposition}

\begin{proof}
Hurwitz stability makes $J|_S$ invertible, hence $J(S)=S$.  Because the full
image of $J$ is contained in $S$, this forces $\operatorname{im}J=S$ and
$\operatorname{rank}J=s$.  If $x\in S\cap\ker J$, invertibility of the
restriction gives $x=0$.  Rank--nullity now gives
$\dim S+\dim\ker J=n$, proving the direct sum.  Both summands are invariant,
and $J$ vanishes on its kernel, which gives the block representation.  The
determinant factorization follows, and its second factor is positive on the
positive real axis because it is the characteristic polynomial of a Hurwitz
matrix.
\end{proof}

Choose any rational basis matrix $B\in\mathbb Q^{n\times s}$ for $S$ and any
rational left inverse $C$ with $CB=I_s$.  Since $JB$ has columns in $S$,
\[
 JB=B(CJB).
\]
Thus $R=CJB$ represents $J|_S$.  If $\widetilde B=BT$ is another basis, then a
compatible left inverse gives
\[
 \widetilde R=T^{-1}RT,
\]
so all reduced Hurwitz tests are basis-independent.

\subsection{Why a nonzero mode is not restricted to \texorpdfstring{$S$}{S}}
For Neumann boundary conditions on a connected bounded domain, every
nonconstant Laplacian eigenfunction $\phi$ is orthogonal to the constant
function, hence
\[
 \int_\Omega\phi(x)\,dx=0.
\]
A linear perturbation of the form $u(x)=a\phi(x)$ therefore satisfies every
integrated linear conservation law $\ell^T\int_\Omega u=0$ for every amplitude
$a\in\mathbb R^n$.  The amplitude space of a nonzero mode is the full species
space, and its matrix is $J-q^2D$.

Restricting that matrix to $S$ is generally impossible.  For example, with
$S=\operatorname{span}(1,1)^T$ and $D=\operatorname{diag}(1,2)$,
\[
 D(1,1)^T=(1,2)^T\notin S.
\]
The homogeneous reaction matrix preserves $S$; unequal species diffusion need
not.  This distinction is essential for conservative networks.
